% ════════════════════════════════════════════════════════════════════
%  tầng b — b1 · b2 · b3, rồi chỗ nối hai tầng
% ════════════════════════════════════════════════════════════════════

\begin{frame}[fragile]{integrated gradients — công cụ của bước b1}
  \framesubtitle{b1 · nền}
  \lead{muốn biết một \emph{núm} ảnh hưởng bao nhiêu, cách thô là vặn nó về 0 rồi
        xem kết quả tụt bao nhiêu. hàng nghìn núm thì phải chạy hàng nghìn lần.}

  \begin{columns}[t]
  \begin{column}{0.49\textwidth}
    \begin{keyb}[title={ig làm khác}]
      vặn \emph{tất cả} núm từ 0 lên mức thật \hl{cùng một lúc}, chia đường đi đó
      làm $M$ chặng, mỗi chặng ghi lại độ nhạy của kết quả theo từng núm rồi cộng.
      \vspace{0.25em}

      \hl{$M$ lần chạy cho ra điểm của \emph{mọi} núm.}
    \end{keyb}
    \[
      \mathrm{ig}_i = (x_i - x_i^{0}) \times
      \frac{1}{M}\sum_{m=1}^{M}
      \frac{\partial F\!\left(x^{0} + \alpha_m (x - x^{0})\right)}{\partial x_i}
    \]
    \dd{$x^0$ = mốc, mọi núm bằng 0 · $\alpha_m = m/M$ là chặng thứ $m$}
  \end{column}
  \begin{column}{0.47\textwidth}
    \begin{bx}[title={tính chất đáng tiền: completeness}]
      \[ \textstyle\sum_i \mathrm{ig}_i = F(x) - F(x^0) \]
      \vspace{-0.5em}

      tổng công của mọi núm cộng lại \yes{đúng bằng} phần kết quả tăng lên —
      không dư, không thiếu. đây là một \hl{đẳng thức kiểm được}, nên cài sai
      là biết ngay.
    \end{bx}
    \begin{bx}[title={vì sao tầng a dùng cách thô, tầng b dùng ig}]
      \begin{dia}
cạnh đồ thị  hàng nghìn, đồ thị CỐ ĐỊNH -> ig
đoạn / pool  prompt PHÌNH theo số đoạn
             -> không có điểm hoà vốn
             -> đành chạy thô
      \end{dia}
    \end{bx}
  \end{column}
  \end{columns}
\end{frame}

\begin{frame}[fragile]{b1 · cái «núm» ở đây là gì}
  \framesubtitle{b1}
  \lead{ta chèn vào mỗi cạnh một hệ số $g \in [0,1]$ — gọi là \hl{gate}. để
        $g \equiv 1$ thì được đúng model gốc; hạ về 0 là bịt cạnh đó lại.}

  \[
    \text{bước truyền tin thứ } l:\qquad
    h_v^{(l)} \;=\; \mathrm{agg}\Bigl(
      \bigl\{\, g_e^{(l)} \cdot m_e^{(l)} \;:\; e = (u \to v) \,\bigr\}
    \Bigr)
  \]

  \vspace{0.4em}
  \begin{columns}[t]
  \begin{column}{0.48\textwidth}
    \begin{badb}[title={một gate cho mỗi cạnh là KHÔNG đủ}]
      cùng một cạnh có thể cần ở bước 1 mà vô dụng ở bước 3. một gate dùng chung
      \no{trộn hai vai} lại, và điểm trả về là một số vô nghĩa.
      \vspace{0.3em}

      \begin{dia}
gate chung      g_e        |E| núm
gate theo bước  g_e^(l)    |E| x L núm
      \end{dia}
      \vspace{0.2em}
      chỉ gate theo bước mới đọc ra được \emph{thứ tự} — tức mới ra được một
      \hl{chuỗi} thay vì một đống.
    \end{badb}
  \end{column}
  \begin{column}{0.48\textwidth}
    \begin{bx}[title={ig trên gate}]
      \[
        \mathrm{ig}_e^{(l)} = g^{*(l)}_e \times
        \frac{1}{M}\sum_{m=1}^{M}
        \frac{\partial F_B\bigl(\alpha_m g^{*}\bigr)}{\partial g_e^{(l)}}
      \]
      \vspace{-0.4em}
      \begin{dia}
g* = 1     khi ĐÀO NHÃN
g* = ĝ     khi BÁO CÁO cho người đọc
      \end{dia}
      \vspace{0.2em}
      \dd{$\hat g$ là hệ số model \emph{tự} học được. dùng nó lúc báo cáo thì
          trung thực; nhưng dùng lúc đào nhãn thì sai — nó là thứ model
          \emph{làm}, không phải thứ ta dùng để \emph{đo}}
    \end{bx}
  \end{column}
  \end{columns}
\end{frame}

\begin{frame}[fragile]{b1 · phép kiểm nội bộ, không tốn gì}
  \framesubtitle{b1}
  \lead{một đường $\pi$ gồm $h$ cạnh và đóng góp $c$ vào kết quả. completeness
        của ig nói thẳng rằng:}
  \[
    \sum_{e \in \pi} \mathrm{ig}_e \;=\; c
    \qquad\text{và nếu các cạnh trong } \pi \text{ đối xứng nhau thì}\qquad
    \mathrm{ig}_e \;=\; \frac{c}{h}
  \]

  \vspace{0.5em}
  \begin{columns}[t]
  \begin{column}{0.48\textwidth}
    \begin{keyb}[title={vì sao đây là thứ quý nhất của tầng b}]
      cài sai bất cứ chỗ nào trong khâu ig — sai dấu, sai số chặng $M$, sai mốc,
      nhân hai lần hệ số — thì tổng \no{không còn bằng} $c$, và ta thấy ngay.
      \vspace{0.3em}

      không có nó thì một lỗi cài đặt sẽ chạy êm, cho ra những đường trông
      hợp lý mà sai.
    \end{keyb}
  \end{column}
  \begin{column}{0.48\textwidth}
    \begin{bx}[title={một đường, ba cạnh}]
      \begin{dia}
c = 0,81          đóng góp của cả đường

ig cạnh 1 = 0,27
ig cạnh 2 = 0,27      mỗi cạnh = c/h
ig cạnh 3 = 0,27           = 0,81/3
            -----
tổng      = 0,81  = c        KHỚP
      \end{dia}
      \vspace{0.2em}
      lệch $\Rightarrow$ có lỗi cài đặt. \no{dừng, đừng đi tiếp.}
    \end{bx}
  \end{column}
  \end{columns}
\end{frame}

\begin{frame}[fragile]{b2 · nối các cạnh thành một đường}
  \framesubtitle{b2}
  \lead{đồ thị gốc có chu trình nên không duyệt một lượt được. mẹo: \emph{vẽ lại}
        nó — mỗi bước truyền tin một bản sao.}

  \begin{columns}[t]
  \begin{column}{0.485\textwidth}
    \begin{bx}[title={trải phẳng}]
      \begin{dia}
node (v, l)  =  node v ở bước l
cạnh (u,l) -> (v,l+1)  nếu u->v trong G

L bước  ->  đồ thị KHÔNG CÒN CHU TRÌNH
        ->  duyệt được một lượt
      \end{dia}
      \vspace{0.2em}
      \yes{thứ tự duyệt chính là số bước $l$} — không cần sắp xếp gì thêm.
    \end{bx}
    \begin{bx}[title={một lượt quét là xong}]
      \[ \mathrm{best}(v, l) = \max_{u \,:\, (u,v) \in E}
         \Bigl[\, \mathrm{best}(u, l-1) + \mathrm{ig}_{(u,v)}^{(l)} \,\Bigr] \]
      \dd{nhớ lại nơi đến từ đâu để truy ngược ra dãy cạnh ·
          chi phí $O\bigl((|V| + |E|)\,L\bigr)$}
    \end{bx}
  \end{column}
  \begin{column}{0.475\textwidth}
    \begin{badb}[title={cộng điểm, KHÔNG phải tìm đường ngắn nhất}]
      quen tay thì đặt $\mathrm{cost} = -\log(\cdot)$ rồi chạy dijkstra.
      \no{sai ở đây.}
      \vspace{0.3em}

      \begin{dia}
ig CỘNG được, theo completeness:
   tổng ig trên một đường
   = đóng góp của đường đó

-log biến TỔNG thành TÍCH
   -> mất completeness
   -> c/h không còn đúng
   -> mất luôn phép kiểm nội bộ
      \end{dia}
      \vspace{0.2em}
      và ig \emph{có dấu} âm cho cạnh gây hại, mà $-\log$ của số âm thì vô nghĩa.
    \end{badb}
  \end{column}
  \end{columns}
\end{frame}

\begin{frame}[fragile]{b3 · một đường là đủ chưa}
  \framesubtitle{b2 · b3}
  \lead{hai câu hỏi khác nhau, và hay bị lẫn: \emph{tập đường đã lấy chiếm bao nhiêu
        phần bằng chứng} và \emph{bỏ đúng đường này đi thì mất bao nhiêu}.}
  \[
    \mathrm{suff}(\Pi_D) \;=\;
    \frac{\sum_{(e,l) \,\in\, \Pi_D} \mathrm{ig}_e^{(l)}}
         {\sum_{(e,l)} \mathrm{ig}_e^{(l)}}
    \qquad\qquad
    \mathrm{nec}(\pi) \;=\; F_B(g^{*}) - F_B\bigl(g^{*} \text{ bỏ } \pi\bigr)
  \]

  \vspace{0.5em}
  \begin{columns}[t]
  \begin{column}{0.47\textwidth}
    \begin{keyb}[title={$\mathrm{suff}$ tính trên TOÀN tập đường}]
      \no{không} cộng $\mathrm{suff}$ của từng đường rồi lấy tổng — hai đường đi
      chung một cạnh thì cạnh đó bị đếm hai lần.
      \vspace{0.3em}

      lấy \hl{hợp các cạnh} rồi mới chia:
      \begin{dia}
hợp cạnh: 0,22 + 0,19 + 0,33 = 0,74
          0,74 / 0,81 = 0,91
suff = 0,91      đủ
      \end{dia}
    \end{keyb}
  \end{column}
  \begin{column}{0.49\textwidth}
    \begin{bx}[title={$\mathrm{suff}$ thấp có HAI nguyên nhân khác hẳn nhau}]
      node nào không đi tới được từ bất kỳ điểm khởi nào vẫn nằm ở \hl{mẫu số}.
      nên $\mathrm{suff}$ thấp chưa chắc do thiếu đường.
      \vspace{0.25em}

      \begin{dia}
suff_max = trần mà mọi đường từ điểm khởi
           có thể đạt — đọc MIỄN PHÍ
           từ bảng quy hoạch động

suff_max thấp   -> nhận diện thực thể hỏng
                   dừng ngay, đừng tìm nữa
suff_max cao    -> đồ thị thiếu cạnh
      \end{dia}
    \end{bx}
  \end{column}
  \end{columns}
\end{frame}

\begin{frame}[fragile]{b2 · khi cần nhiều hơn một đường}
  \framesubtitle{b2}
  \lead{một đoạn có thể được chứng minh bởi vài đường hội tụ lại. gom dần cho tới
        khi đủ, rồi tỉa ngược những đường hoá ra thừa.}

  \begin{columns}[t]
  \begin{column}{0.52\textwidth}
    \begin{bx}[title={hai chặng}]
      \begin{dia}
-- chặng 1: gom -------------------------
tập đường <- rỗng
lặp:
    lấy đường điểm cao nhất từ bảng qhđ
    gain = tổng ig các cạnh MỚI của nó
    thêm vào tập đường

    đủ (suff >= nguong)  -> DỪNG, thành công
    gain quá nhỏ         -> DỪNG, bằng chứng
                            phân tán
    quá n_max đường      -> DỪNG, hết ngân sách
                            người đọc chịu được

    đặt ig = 0 cho mọi cạnh vừa lấy
    chạy lại quy hoạch động

-- chặng 2: tỉa ngược -------------------
với từng đường, theo thứ tự đã thêm:
    bỏ nó ra mà vẫn đủ  ->  bỏ hẳn
      \end{dia}
    \end{bx}
  \end{column}
  \begin{column}{0.44\textwidth}
    \begin{keyb}[title={cạnh dùng chung tự thành miễn phí}]
      đặt $\mathrm{ig} = 0$ cho cạnh đã lấy, nên vòng sau \yes{không trả tiền lại}
      cho cạnh đó. hội tụ nhiều đường được xử lý \emph{tự nhiên}, không cần luật
      riêng nào.
    \end{keyb}
    \begin{bx}[title={$n_{\max}$ không phải con số chọn bừa}]
      \[ n_{\max} = \Bigl\lceil \frac{\theta_s}{\tilde\varepsilon} \Bigr\rceil \]
      \dd{mỗi đường được nhận phải góp ít nhất $\tilde\varepsilon$, nên quá
          $\theta_s/\tilde\varepsilon$ đường là không thể còn đủ. con số rơi ra
          từ hai ngưỡng đã có, không phải đặt tay}
    \end{bx}
  \end{column}
  \end{columns}
\end{frame}

% ════════════════════════════════════════════════════════════════════
\section{nối hai tầng}

\begin{frame}[fragile]{chỉ một thứ đi qua ranh giới}
  \lead{tầng a tính ra rất nhiều số. \hl{không số nào} được đi sang tầng b. lý do
        không phải cho gọn — mà vì chúng \emph{đo trên hai không gian khác nhau}.}

  \begin{columns}[t]
  \begin{column}{0.415\textwidth}
    \begin{keyb}[title={giao diện}]
      \begin{dia}
TẦNG A -> T = {đoạn 3, đoạn 12} -> TẦNG B

KHÔNG qua: v(S) · tau_q · beta · pgi
           g_theta · W · entropy
      \end{dia}
      \vspace{0.25em}
      điểm của tầng a đo trên \emph{rổ 40 đoạn}; tầng b đi trên \emph{đồ thị}.
      mang một thang đo từ nơi này sang nơi kia là so hai thứ không so được.
    \end{keyb}
  \end{column}
  \begin{column}{0.545\textwidth}
    \begin{bx}[title={$T$ dùng làm nhãn để dạy bộ đi tìm}]
      \begin{align*}
        \mathcal{L}_{\text{node}} &= \mathrm{bce}\bigl(p(v),\; \mathbf{1}[v \in T]\bigr)\\[0.2em]
        \mathcal{L}_{\text{hop}}  &= \sum_{l}\;\sum_{v \,\in\, V_{\text{hop}}(l)}
          \mathrm{bce}\bigl(p^{(l)}(v),\; \mathbf{1}[(v,l) \in \Pi]\bigr)\\[0.2em]
        \mathcal{L}_{\text{edge}} &= \sum_{(e,l) \,\in\, \Pi}
          -\log \hat g_e^{(l)}
      \end{align*}
      \vspace{-0.3em}
      \begin{dia}
âm KHÓ = đúng node, SAI bước
      \end{dia}
      \dd{không có loại âm khó này thì tỉ lệ dương:âm là \no{$1 : 10^6$} và model
          học được cách luôn trả lời «không»}
    \end{bx}
  \end{column}
  \end{columns}
\end{frame}

\begin{frame}{bảng chẩn đoán — và nó cũng là nguồn nhãn}
  \lead{hai trục độc lập: \emph{đường có giải thích được không} và \emph{bộ đi tìm
        có xếp đoạn đúng lên cao không}. bốn ô, bốn bệnh khác nhau.}

  \begin{center}
  \begin{tabular}[t]{@{}>{\raggedright\arraybackslash}p{0.13\linewidth}
                        >{\raggedright\arraybackslash}p{0.345\linewidth}
                        >{\raggedright\arraybackslash}p{0.345\linewidth}@{}}
    \toprule
    & \textbf{xếp hạng tốt} & \textbf{xếp hạng kém}\\
    \midrule
    \hl{$\mathrm{suff}$ cao}
      & \yes{lành} — tìm đúng, và giải thích được
      & \no{lỗi model} — đường có đó mà không xếp đoạn lên cao\\[0.6em]
    \hl{$\mathrm{suff}$ thấp}
      & \no{đường tắt} — xếp đúng \emph{nhờ lối khác}, không nhờ bằng chứng
      & \no{lỗi đồ thị} — thiếu cạnh, hoặc nhận diện thực thể hỏng\\
    \bottomrule
  \end{tabular}
  \end{center}

  \vspace{0.8em}
  \begin{keyb}
    ô \emph{đường tắt} là ô dễ bỏ sót nhất: mọi chỉ số đều đẹp mà lời giải thích
    rỗng. và mỗi ô \yes{tự nó là một lô nhãn} để huấn luyện vòng sau.
  \end{keyb}
\end{frame}
